\section{Consequences and the contraction boundary}

\begin{remark}[Quenched random environments]
If the rooted tree and its channel parameters are random but, conditional on
the environment, form a locally finite deterministic tree satisfying a common
bound $\lambda_e\le q<1$, Theorem~\ref{thm:main} applies to almost every
realization. Thus quenched reconstruction is equivalent realization by
realization to positive quenched capacity. This observation does not by itself
identify an annealed threshold, because averaging may mix environments with
different capacities.
\end{remark}

\begin{corollary}[Spherical trees]\label{cor:spherical}
Suppose every vertex at the same depth has the same number of children,
let $N_k=|L_k|$, and suppose every depth-$k$ edge has parameter
$\lambda_k\le q<1$. Put $\rho_k=\prod_{j=1}^k\lambda_j^2$.
Then
\[
\Cap_3(T,\rho)=
\left(\sum_{k\ge1}(N_k\rho_k)^{-2}\right)^{-1/2},
\]
where the value is zero if the series diverges. Consequently reconstruction
occurs exactly when that series converges.
\end{corollary}

\begin{proof}
The uniform flow of mass $a$ assigns $a/N_k$ to each depth-$k$ edge,
giving the displayed path energy. Conversely, choose a uniform ray,
so its vertex at level $k$ is uniform on $L_k$. For any flow of mass
$a$, Jensen gives
\[
\E\left(\frac{\theta(e_{V_k})}{\rho_k}\right)^2
\ge\left(\frac{a}{N_k\rho_k}\right)^2.
\]
Average the path potential and sum by monotone convergence. Its supremum
is at least this average, proving optimality of the uniform flow.
\end{proof}

In particular, if $N_k\rho_k\asymp k^\beta$, the criterion is
$\beta>1/2$. On a regular $d$-ary tree with common parameter $\lambda$,
it reduces to $d\lambda^2>1$, including nonreconstruction at equality.
These are direct consequences of the proved series criterion.

\begin{proposition}[The weighted capacity is exact on one ray]
\label{prop:ray-weighted}
Suppose $T$ is one infinite ray, with edge parameters
$0<\lambda_k\le1$. Put
\[
 \Lambda_n=\prod_{k=1}^n\lambda_k,
 \qquad \rho_n=\Lambda_n^2,
 \qquad \Lambda_\infty=\lim_n\Lambda_n.
\]
Then
\[
 a_\infty=\frac34\Lambda_\infty,
\]
and
\begin{equation}\label{eq:ray-weighted-capacity}
 \Cap_3^{\mathrm{noise}}(T,\rho)=
 \begin{cases}
  (\rho_\infty^{-2}-1)^{-1/2},&0<\rho_\infty<1,\\
  +\infty,&\rho_\infty=1,\\
  0,&\rho_\infty=0.
 \end{cases}
\end{equation}
Consequently, reconstruction on a ray is equivalent to positive weighted
capacity without any uniform contraction assumption.
\end{proposition}

\begin{proof}
Composition of the symmetric channels multiplies their parameters, so the
root-to-level-$n$ channel has parameter $\Lambda_n$. Observing the terminal
spin therefore gives total-variation advantage $3\Lambda_n/4$.

Every flow on a ray is constant, say of mass $a$. Since
$\rho_k=\rho_{k-1}\lambda_k^2$,
\[
 \frac{1-\lambda_k^4}{\rho_k^2}
 =\rho_k^{-2}-\rho_{k-1}^{-2}.
\]
Thus its weighted path energy through level $n$ is
$a^2(\rho_n^{-2}-1)$. Letting $n\to\infty$ proves
\eqref{eq:ray-weighted-capacity}.
\end{proof}

The uniform upper bound in Theorem~\ref{thm:main} cannot simply be discarded
while retaining the unweighted capacity. For example, on one ray take
$\lambda_k=\exp(-2^{-k})$. Then $\Lambda_\infty=e^{-1}$, so
$a_\infty=3/(4e)>0$. Every nonzero flow has constant mass, whereas each
unweighted capacity summand is at least the square of that mass; hence
$\Cap_3(T,\rho)=0$. Proposition~\ref{prop:ray-weighted} shows that the
edge-noise weight repairs this obstruction exactly on a ray. Theorem
\ref{thm:weighted-necessity} proves the corresponding necessity on every
tree. A general weighted-capacity converse would require a refinement of the
sufficiency loss estimate: the geometric-kernel bound used above deliberately
collapses long one-child stretches and is not controlled by weighted energy
alone.

\begin{remark}[Zero edges]
The statement assumes strictly positive edge parameters so that
$\rho^{-1}$ is defined. An edge with $\lambda_e=0$ separates the spin
below it completely from the root. If zero edges are to be admitted,
remove their downstream components and use the theorem on the remaining
root component; the capacity formula must be interpreted on that
component, not by dividing by zero. If the remaining root component is
finite, reconstruction is impossible and its capacity is zero.
\end{remark}
